\chapter{Linear Algebra and Algorithms}

\section{Numerical Solution of Linear Systems of Equations}

The problem of solving the linear system $\mathbf{A}\mathbf{x} = \mathbf{b}$ for a given matrix $\mathbf{A} \in \R^{n \times n}$ and vector $\mathbf{b} \in \R^n$ is the ultimate foundational problem in computational mathematics. In the realm of large-scale data science, and particularly for systems arising from the discretization of partial differential equations (PDEs) or high-dimensional optimisation, the direct methods that yield solutions in a finite number of arithmetic operations are supplanted by the elegance of Iterative Methods. These methods trade finite termination for asymptotic convergence, leveraging the sparsity and spectral properties inherent in many real-world matrices.

\subsection{Iterative Methods: Jacobi and Gauss-Seidel}

Iterative methods are fundamentally rooted in the theory of fixed-point iteration on a Banach space. We seek a solution $\mathbf{x} \in \R^n$ as the unique fixed point of a contracted mapping.

\begin{definition}[Matrix Splitting and First-Order Scheme]
Let $\mathbf{A} \in \R^{n \times n}$ be an invertible matrix. An iterative method for solving $\mathbf{A}\mathbf{x} = \mathbf{b}$ is generated by a splitting of $\mathbf{A}$ into two matrices,
\[ \mathbf{A} = \mathbf{M} - \mathbf{N} \]
where $\mathbf{M}$ is a non-singular matrix that is computationally trivial to invert. The recurrence relation is:
\[ \mathbf{x}^{(k+1)} = \mathbf{M}^{-1}\mathbf{N}\mathbf{x}^{(k)} + \mathbf{M}^{-1}\mathbf{b} = \mathbf{T}\mathbf{x}^{(k)} + \mathbf{c} \]
The matrix $\mathbf{T} = \mathbf{M}^{-1}\mathbf{N}$ is the iteration matrix.
\end{definition}

For $\mathbf{A}$, we adopt the canonical additive decomposition: $\mathbf{A} = \mathbf{D} - \mathbf{L} - \mathbf{U}$.

\subsubsection{The Jacobi Iteration}
\begin{definition}[Jacobi Splitting]
\[ \mathbf{M}_J = \mathbf{D}, \quad \mathbf{N}_J = \mathbf{L} + \mathbf{U} \]
The iteration matrix is $\mathbf{T}_J = \mathbf{D}^{-1}(\mathbf{L} + \mathbf{U})$.
\end{definition}

\subsubsection{The Gauss-Seidel Iteration}
\begin{definition}[Gauss-Seidel Splitting]
\[ \mathbf{M}_{GS} = \mathbf{D} - \mathbf{L}, \quad \mathbf{N}_{GS} = \mathbf{U} \]
The iteration matrix is $\mathbf{T}_{GS} = (\mathbf{D} - \mathbf{L})^{-1}\mathbf{U}$.
\end{definition}

\subsubsection{The Rigor of Convergence: The Spectral Radius Criterion}

\begin{theorem}[Fundamental Criterion for Iterative Convergence]
The iterative method converges for every initial vector $\mathbf{x}^{(0)}$ if and only if the spectral radius of the iteration matrix $\mathbf{T}$ satisfies:
\[ \rho(\mathbf{T}) = \max_{\lambda \in \sigma(\mathbf{T})} |\lambda| < 1 \]
The asymptotic rate of convergence is $\mathcal{R}(\mathbf{T}) = -\log \rho(\mathbf{T})$.
\end{theorem}

\begin{theorem}[Convergence via SDD]
If $\mathbf{A}$ is strictly diagonally dominant (SDD), then both $\mathbf{T}_J$ and $\mathbf{T}_{GS}$ converge.
\end{theorem}

\begin{lemma}[Spectral Radius Bound by Induced Norm]
This is trivial but needed, let $\mathbf{T} \in \mathbb{C}^{n \times n}$ be an arbitrary square matrix. For every induced matrix norm, denoted $\|\cdot\|$, the spectral radius of $\mathbf{T}$ is bounded above by the value of that norm:
\[ \rho(\mathbf{T}) \leq \|\mathbf{T}\| \]
\end{lemma}

\begin{remark}
Strict Diagonal Dominance of $\mathbf{A}$ is a rigorous sufficient condition for the iterative stability of the matrix splitting. It translates mathematically into the statement that the $\ell_1$-norm of the strictly off-diagonal part of $\mathbf{A}$, when normalized by the diagonal entry, is globally bounded by a factor strictly less than unity. Specifically, defining $\mathbf{T}_J$ such that $\mathbf{A} = \mathbf{D}(\mathbf{I} - \mathbf{T}_J)$, the SDD condition guarantees that the induced $\ell_{\infty}$ matrix norm of the Jacobi iteration matrix is a proper contraction:
\[ \|\mathbf{T}_J\|_{\infty} = \max_{i} \sum_{j \ne i} \left| \frac{a_{ij}}{a_{ii}} \right| < 1 \]
Since $\rho(\mathbf{T}_J) \le \|\mathbf{T}_J\|_{\infty}$, this immediately implies $\rho(\mathbf{T}_J) < 1$. For Gauss-Seidel, the condition is subtler, relating to a non-induced norm bound on the sequential error ratio (the factor $C_i$), which is also forced below unity by SDD.
\end{remark}

\begin{proof}[Proof Sketch for Gauss-Seidel]
By defining the contraction factor $C_{i^{\ast}}$ in the $\ell_{\infty}$ norm (as derived previously), we show that $\rho(\mathbf{T}_{GS}) \le \|\mathbf{T}_{GS}\|_{\infty} \le \max_i C_i < 1$, with:
\[ C_i = \frac{\sum_{j > i} |a_{ij}|}{|a_{ii}| - \sum_{j < i} |a_{ij}|} \]
The condition SDD ensures $C_i < 1$.
\end{proof}

\subsection{Architectural Analysis: Computational Cost and Parallelism}

Beyond mere existence of convergence, the utility of an iterative method in a modern computing environment is governed by its computational complexity and data dependency.

\subsubsection{Computational Complexity (Cost per Iteration)}

The expense of a single iteration is dominated by the matrix-vector multiplication associated with $\mathbf{N}\mathbf{x}^{(k)}$ and the inversion (application of the inverse) of $\mathbf{M}$.

\begin{proposition}[Cost per Iteration]
For a dense matrix $\mathbf{A} \in \R^{n \times n}$, the complexity of a single iteration for both Jacobi and Gauss-Seidel is $\mathcal{O}(n^2)$.
\end{proposition}
\begin{proof}
Both methods primarily involve computing the product $\mathbf{N}\mathbf{x}^{(k)}$. Since $\mathbf{N}$ is full (of size $n \times n$), this requires $n^2$ multiplications and $n^2 - n$ additions, leading to a complexity of $\mathcal{O}(n^2)$. The inversion of $\mathbf{M}$ (diagonal for Jacobi, triangular for Gauss-Seidel) is only $\mathcal{O}(n)$ via vector scaling or forward substitution, respectively.

However, in Data Science and Numerical PDE's, $\mathbf{A}$ is often sparse. If $\mathbf{A}$ has $\mathfrak{n}_z$ non-zero entries (where $\mathfrak{n}_z \ll n^2$), the cost of computing $\mathbf{N}\mathbf{x}^{(k)}$ is reduced to $\mathcal{O}(\mathfrak{n}_z)$, which is often $\mathcal{O}(n)$ for structured matrices (e.g., tridiagonal matrices from 1D discretizations, where $\mathfrak{n}_z \approx 3n$).
\end{proof}

\subsubsection{Data Dependency and Parallelization}

The true dichotomy between Jacobi and Gauss-Seidel emerges in their temporal data dependency, which dictates their amenability to parallel computing.

\begin{enumerate}
    \item \textbf{Jacobi Method (Concurrent):} The $i$-th component of the new iterate, $x_i^{(k+1)}$, is computed using only the components from the previous iteration, $\mathbf{x}^{(k)}$.
    \[ x_i^{(k+1)} = \frac{1}{a_{ii}} \left( b_i - \sum_{j \ne i} a_{ij} x_j^{(k)} \right) \]
    Crucially, all $x_i^{(k+1)}$ can be computed simultaneously and independently. This makes the Jacobi method trivially parallelizable, an immense advantage in modern GPU/multicore architectures.
    
    \item \textbf{Gauss-Seidel Method (Sequential):} The $i$-th component of the new iterate, $x_i^{(k+1)}$, utilizes a mixture: the newly computed components $x_j^{(k+1)}$ for $j < i$ and the old components $x_j^{(k)}$ for $j > i$.
    \[ x_i^{(k+1)} = \frac{1}{a_{ii}} \left( b_i - \sum_{j < i} a_{ij} x_j^{(k+1)} - \sum_{j > i} a_{ij} x_j^{(k)} \right) \]
    This explicit dependency of $x_i^{(k+1)}$ on $x_1^{(k+1)}, \dots, x_{i-1}^{(k+1)}$ creates a sequential bottleneck (a triangular solve), making the standard Gauss-Seidel iteration difficult to parallelize for general matrices $\mathbf{A}$. Only specialized reordering (e.g., Red-Black ordering for grid problems) can restore parallelism.
\end{enumerate}

\subsection{Heuristics and Advanced Intuition}

\subsubsection{The Causal Advantage of Gauss-Seidel}

The reason $\rho(\mathbf{T}_{GS}) \le \rho(\mathbf{T}_J)$ often holds (and thus, Gauss-Seidel converges faster) is rooted in its immediate information propagation. The change made in the $i$-th component is instantly used to refine the $(i+1)$-th component within the same time step $k \to k+1$. Jacobi delays this information until the next time step, $k+1 \to k+2$.

\subsubsection{Symmetric Positive Definite Matrices}

\begin{definition}[Symmetric Positive Definite (SPD)]
A matrix $\mathbf{A} \in \R^{n \times n}$ is Symmetric Positive Definite if $\mathbf{A} = \mathbf{A}^T$ and for all non-zero vectors $\mathbf{x} \in \R^n$, $\mathbf{x}^T \mathbf{A} \mathbf{x} > 0$.
\end{definition}

\begin{theorem}[Convergence for SPD Matrices]
If $\mathbf{A}$ is Symmetric Positive Definite, the Gauss-Seidel iteration converges. (Note: Jacobi convergence requires an additional condition, such as $\mathbf{D} - (\mathbf{L}+\mathbf{U})$ being positive definite, which is often harder to verify than SDD.)
\end{theorem}
\begin{remark}[Minimisation View]
For SPD matrices, solving $\mathbf{A}\mathbf{x} = \mathbf{b}$ is equivalent to finding the unique minimiser of the quadratic function:
\[ f(\mathbf{x}) = \frac{1}{2}\mathbf{x}^T \mathbf{A} \mathbf{x} - \mathbf{b}^T \mathbf{x} \]
The Gauss-Seidel method, in this context, is equivalent to Coordinate Descent, where $x_i$ is updated to minimise $f(\mathbf{x})$ while all other components $\mathbf{x}_j$ ($j \ne i$) are held fixed. Since $f(\mathbf{x})$ is strictly convex, this sequence of exact line searches is guaranteed to converge to the unique global minimiser $\mathbf{x}$.
\end{remark}

\subsubsection{The Role of Ordering and Scaling}

\begin{itemize}
    \item \textbf{Reordering (Permutation):} The iteration matrices $\mathbf{T}_{GS}$ and $\mathbf{T}_J$ depend explicitly on the matrix $\mathbf{A}$. If we permute the rows and columns via a permutation matrix $\mathbf{P}$ such that the system becomes $(\mathbf{P}\mathbf{A}\mathbf{P}^T)(\mathbf{P}\mathbf{x}) = \mathbf{P}\mathbf{b}$, the new matrix $\mathbf{A}' = \mathbf{P}\mathbf{A}\mathbf{P}^T$ can have drastically different convergence properties. Reordering to maximize the concentration of large coefficients on the diagonal (to approach SDD) can be crucial.
    \item \textbf{Preconditioning:} The core idea of Preconditioning is to replace $\mathbf{A}\mathbf{x} = \mathbf{b}$ with an equivalent system $\mathbf{P}^{-1}\mathbf{A}\mathbf{x} = \mathbf{P}^{-1}\mathbf{b}$ where $\mathbf{P}^{-1}\mathbf{A}$ has a much smaller condition number $\kappa(\mathbf{P}^{-1}\mathbf{A})$ and, critically, an iteration matrix $\mathbf{T}'$ with a smaller spectral radius $\rho(\mathbf{T}')$. The iterative methods discussed here are the simplest forms of preconditioners for themselves (e.g., Jacobi is the Diagonal Preconditioner).
\end{itemize}

\subsection{Numerical Diagnosis: Implementation and Error Analysis}

The final step is the rigorous implementation, which requires a precise definition of stopping criteria to ensure the computed solution $\mathbf{x}^{(k)}$ is an acceptable approximation of the true solution $\mathbf{x}$.

\subsubsection{Numerical Implementation in Python}

The true essence of iterative methods manifests in their computational implementation. We employ the \texttt{numpy} library for vectorized operations, acknowledging the fundamental trade-off: \textbf{Jacobi} benefits immensely from vectorization (parallelism), while \textbf{Gauss-Seidel} fundamentally requires sequential processing (causality).

\paragraph{Jacobi Iteration: Vectorized (Trivially Parallel)}
The Jacobi method is ideally suited for vectorization, as the calculation of $\mathbf{x}^{(k+1)}$ depends only on $\mathbf{x}^{(k)}$, allowing simultaneous computation of all components. This formulation directly expresses the matrix splitting $\mathbf{D}\mathbf{x}^{(k+1)} = (\mathbf{L}+\mathbf{U})\mathbf{x}^{(k)} + \mathbf{b}$.

\begin{Verbatim}[commandchars=\\\{\}]
\PY{k+kn}{import}\PY{+w}{ }\PY{n+nn}{numpy}\PY{+w}{ }\PY{k}{as}\PY{+w}{ }\PY{n+nn}{np}

\PY{k}{def}\PY{+w}{ }\PY{n+nf}{jacobi\PYZus{}solve}\PY{p}{(}\PY{n}{A}\PY{p}{,} \PY{n}{b}\PY{p}{,} \PY{n}{x0}\PY{p}{,} \PY{n}{max\PYZus{}iter}\PY{p}{,} \PY{n}{tol}\PY{o}{=}\PY{l+m+mf}{1e\PYZhy{}6}\PY{p}{)}\PY{p}{:}
\PY{+w}{    }\PY{l+s+sd}{\PYZdq{}\PYZdq{}\PYZdq{}}
\PY{l+s+sd}{    Solves Ax = b using the Jacobi method. M = D, N = L + U.}
\PY{l+s+sd}{    Implementation is highly vectorized (O(n\PYZca{}2) or O(n) for sparse A).}
\PY{l+s+sd}{    \PYZdq{}\PYZdq{}\PYZdq{}}
    \PY{n}{n} \PY{o}{=} \PY{n}{A}\PY{o}{.}\PY{n}{shape}\PY{p}{[}\PY{l+m+mi}{0}\PY{p}{]}
    \PY{n}{x\PYZus{}old} \PY{o}{=} \PY{n}{np}\PY{o}{.}\PY{n}{array}\PY{p}{(}\PY{n}{x0}\PY{p}{,} \PY{n}{dtype}\PY{o}{=}\PY{n}{A}\PY{o}{.}\PY{n}{dtype}\PY{p}{)} \PY{c+c1}{\PYZsh{} Initialize iterate}
    
    \PY{c+c1}{\PYZsh{} Extract D and R = L + U (R is off\PYZhy{}diagonal)}
    \PY{n}{D} \PY{o}{=} \PY{n}{np}\PY{o}{.}\PY{n}{diag}\PY{p}{(}\PY{n}{A}\PY{p}{)}
    \PY{n}{R} \PY{o}{=} \PY{n}{A} \PY{o}{\PYZhy{}} \PY{n}{np}\PY{o}{.}\PY{n}{diag}\PY{p}{(}\PY{n}{D}\PY{p}{)}
    
    \PY{k}{for} \PY{n}{k} \PY{o+ow}{in} \PY{n+nb}{range}\PY{p}{(}\PY{l+m+mi}{1}\PY{p}{,} \PY{n}{max\PYZus{}iter} \PY{o}{+} \PY{l+m+mi}{1}\PY{p}{)}\PY{p}{:}
        \PY{c+c1}{\PYZsh{} x\PYZus{}new = D\PYZca{}\PYZob{}\PYZhy{}1\PYZcb{} * (b \PYZhy{} R * x\PYZus{}old)}
        \PY{c+c1}{\PYZsh{} @ is matrix multiplication; / D is component\PYZhy{}wise division (D\PYZca{}\PYZob{}\PYZhy{}1\PYZcb{} operation)}
        \PY{n}{x\PYZus{}new} \PY{o}{=} \PY{p}{(}\PY{n}{b} \PY{o}{\PYZhy{}} \PY{n}{R} \PY{o}{@} \PY{n}{x\PYZus{}old}\PY{p}{)} \PY{o}{/} \PY{n}{D} 
        
        \PY{c+c1}{\PYZsh{} Relative Error Criterion (L\PYZus{}inf norm)}
        \PY{n}{error\PYZus{}norm} \PY{o}{=} \PY{n}{np}\PY{o}{.}\PY{n}{linalg}\PY{o}{.}\PY{n}{norm}\PY{p}{(}\PY{n}{x\PYZus{}new} \PY{o}{\PYZhy{}} \PY{n}{x\PYZus{}old}\PY{p}{,} \PY{n}{np}\PY{o}{.}\PY{n}{inf}\PY{p}{)}
        \PY{k}{if} \PY{n}{error\PYZus{}norm} \PY{o}{\PYZlt{}} \PY{n}{tol} \PY{o}{*} \PY{n}{np}\PY{o}{.}\PY{n}{linalg}\PY{o}{.}\PY{n}{norm}\PY{p}{(}\PY{n}{x\PYZus{}new}\PY{p}{,} \PY{n}{np}\PY{o}{.}\PY{n}{inf}\PY{p}{)}\PY{p}{:}
            \PY{c+c1}{\PYZsh{} print(f\PYZdq{}Jacobi converged in \PYZob{}k\PYZcb{} iterations.\PYZdq{})}
            \PY{k}{return} \PY{n}{x\PYZus{}new}\PY{p}{,} \PY{n}{k}
            
        \PY{n}{x\PYZus{}old} \PY{o}{=} \PY{n}{x\PYZus{}new}
        
    \PY{c+c1}{\PYZsh{} print(f\PYZdq{}Jacobi failed to converge in \PYZob{}max\PYZus{}iter\PYZcb{} iterations.\PYZdq{})}
    \PY{k}{return} \PY{n}{x\PYZus{}old}\PY{p}{,} \PY{n}{max\PYZus{}iter}
\end{Verbatim}


\paragraph{Gauss-Seidel Iteration: Explicit Sequential Loop (Causality)}
The Gauss-Seidel method \textbf{cannot} be fully vectorized in its canonical form because each $x_i^{(k+1)}$ depends causally on $x_j^{(k+1)}$ for $j < i$. This necessitates the use of a sequential loop over the components $i = 1, \dots, n$, explicitly revealing the data dependency constraint (the lower-triangular solve).

\begin{Verbatim}[commandchars=\\\{\}]
\PY{k}{def}\PY{+w}{ }\PY{n+nf}{gauss\PYZus{}seidel\PYZus{}solve}\PY{p}{(}\PY{n}{A}\PY{p}{,} \PY{n}{b}\PY{p}{,} \PY{n}{x0}\PY{p}{,} \PY{n}{max\PYZus{}iter}\PY{p}{,} \PY{n}{tol}\PY{o}{=}\PY{l+m+mf}{1e\PYZhy{}6}\PY{p}{)}\PY{p}{:}
\PY{+w}{    }\PY{l+s+sd}{\PYZdq{}\PYZdq{}\PYZdq{}}
\PY{l+s+sd}{    Solves Ax = b using the Gauss\PYZhy{}Seidel method. M = D \PYZhy{} L, N = U.}
\PY{l+s+sd}{    Requires an explicit sequential loop (forward substitution).}
\PY{l+s+sd}{    \PYZdq{}\PYZdq{}\PYZdq{}}
    \PY{n}{n} \PY{o}{=} \PY{n}{A}\PY{o}{.}\PY{n}{shape}\PY{p}{[}\PY{l+m+mi}{0}\PY{p}{]}
    \PY{n}{x\PYZus{}old} \PY{o}{=} \PY{n}{np}\PY{o}{.}\PY{n}{array}\PY{p}{(}\PY{n}{x0}\PY{p}{,} \PY{n}{dtype}\PY{o}{=}\PY{n}{A}\PY{o}{.}\PY{n}{dtype}\PY{p}{)}
    
    \PY{k}{for} \PY{n}{k} \PY{o+ow}{in} \PY{n+nb}{range}\PY{p}{(}\PY{l+m+mi}{1}\PY{p}{,} \PY{n}{max\PYZus{}iter} \PY{o}{+} \PY{l+m+mi}{1}\PY{p}{)}\PY{p}{:}
        \PY{n}{x\PYZus{}new} \PY{o}{=} \PY{n}{np}\PY{o}{.}\PY{n}{copy}\PY{p}{(}\PY{n}{x\PYZus{}old}\PY{p}{)}
        
        \PY{c+c1}{\PYZsh{} Sequential update: The causal heart of the method}
        \PY{k}{for} \PY{n}{i} \PY{o+ow}{in} \PY{n+nb}{range}\PY{p}{(}\PY{n}{n}\PY{p}{)}\PY{p}{:}
            \PY{c+c1}{\PYZsh{} 1. Calculate sum for j \PYZlt{} i (L*x\PYZus{}new: using updated values)}
            \PY{n}{sum\PYZus{}L} \PY{o}{=} \PY{n}{np}\PY{o}{.}\PY{n}{dot}\PY{p}{(}\PY{n}{A}\PY{p}{[}\PY{n}{i}\PY{p}{,} \PY{p}{:}\PY{n}{i}\PY{p}{]}\PY{p}{,} \PY{n}{x\PYZus{}new}\PY{p}{[}\PY{p}{:}\PY{n}{i}\PY{p}{]}\PY{p}{)}
            \PY{c+c1}{\PYZsh{} 2. Calculate sum for j \PYZgt{} i (U*x\PYZus{}old: using old values)}
            \PY{n}{sum\PYZus{}U} \PY{o}{=} \PY{n}{np}\PY{o}{.}\PY{n}{dot}\PY{p}{(}\PY{n}{A}\PY{p}{[}\PY{n}{i}\PY{p}{,} \PY{n}{i}\PY{o}{+}\PY{l+m+mi}{1}\PY{p}{:}\PY{p}{]}\PY{p}{,} \PY{n}{x\PYZus{}old}\PY{p}{[}\PY{n}{i}\PY{o}{+}\PY{l+m+mi}{1}\PY{p}{:}\PY{p}{]}\PY{p}{)}
            
            \PY{c+c1}{\PYZsh{} Update: x\PYZus{}i\PYZca{}(k+1) = (b\PYZus{}i \PYZhy{} sum\PYZus{}L \PYZhy{} sum\PYZus{}U) / A\PYZus{}ii}
            \PY{n}{x\PYZus{}new}\PY{p}{[}\PY{n}{i}\PY{p}{]} \PY{o}{=} \PY{p}{(}\PY{n}{b}\PY{p}{[}\PY{n}{i}\PY{p}{]} \PY{o}{\PYZhy{}} \PY{n}{sum\PYZus{}L} \PY{o}{\PYZhy{}} \PY{n}{sum\PYZus{}U}\PY{p}{)} \PY{o}{/} \PY{n}{A}\PY{p}{[}\PY{n}{i}\PY{p}{,} \PY{n}{i}\PY{p}{]}
        
        \PY{c+c1}{\PYZsh{} Relative Error Criterion (L\PYZus{}inf norm)}
        \PY{n}{error\PYZus{}norm} \PY{o}{=} \PY{n}{np}\PY{o}{.}\PY{n}{linalg}\PY{o}{.}\PY{n}{norm}\PY{p}{(}\PY{n}{x\PYZus{}new} \PY{o}{\PYZhy{}} \PY{n}{x\PYZus{}old}\PY{p}{,} \PY{n}{np}\PY{o}{.}\PY{n}{inf}\PY{p}{)}
        \PY{k}{if} \PY{n}{error\PYZus{}norm} \PY{o}{\PYZlt{}} \PY{n}{tol} \PY{o}{*} \PY{n}{np}\PY{o}{.}\PY{n}{linalg}\PY{o}{.}\PY{n}{norm}\PY{p}{(}\PY{n}{x\PYZus{}new}\PY{p}{,} \PY{n}{np}\PY{o}{.}\PY{n}{inf}\PY{p}{)}\PY{p}{:}
            \PY{c+c1}{\PYZsh{} print(f\PYZdq{}Gauss\PYZhy{}Seidel converged in \PYZob{}k\PYZcb{} iterations.\PYZdq{})}
            \PY{k}{return} \PY{n}{x\PYZus{}new}\PY{p}{,} \PY{n}{k}
            
        \PY{n}{x\PYZus{}old} \PY{o}{=} \PY{n}{x\PYZus{}new}
        
    \PY{c+c1}{\PYZsh{} print(f\PYZdq{}Gauss\PYZhy{}Seidel failed to converge in \PYZob{}max\PYZus{}iter\PYZcb{} iterations.\PYZdq{})}
    \PY{k}{return} \PY{n}{x\PYZus{}old}\PY{p}{,} \PY{n}{max\PYZus{}iter}
\end{Verbatim}


\subsubsection{Stopping Criteria and Error Norms}

The convergence condition is $\rho(\mathbf{T}) < 1$, which means $\|\mathbf{e}^{(k)}\| \to 0$. In practice, we halt the process when the error is sufficiently small.

\begin{definition}[Stopping Criteria]
We use a relative error or residual criterion:
\begin{enumerate}
    \item \textbf{Relative Error Criterion (Practical):} Stop when the difference between successive iterates is small:
    \[ \|\mathbf{x}^{(k+1)} - \mathbf{x}^{(k)}\| < \mathrm{tol} \cdot \|\mathbf{x}^{(k+1)}\| \]
    where $\|\cdot\|$ is typically the $\ell_2$ or $\ell_{\infty}$ norm.
    \item \textbf{Residual Criterion (More Robust):} Stop when the residual vector $\mathbf{r}^{(k)} = \mathbf{b} - \mathbf{A}\mathbf{x}^{(k)}$ is small:
    \[ \frac{\|\mathbf{r}^{(k)}\|}{\|\mathbf{b}\|} < \mathrm{tol} \]
    This is often preferred as it directly relates to how well the system $\mathbf{A}\mathbf{x}=\mathbf{b}$ is satisfied.
\end{enumerate}
The choice of the matrix norm (e.g., $\ell_{\infty}$ norm in the SDD proof) is usually dictated by analytical convenience, but the choice of the vector norm for the stopping criterion is a practical trade-off, with $\ell_2$ and $\ell_{\infty}$ being most common.
\end{definition}